% Laurent Warlouzet: From The Hague summit to the Maastricht Treaty – Reading Summary
% Introduction to the European Union, Week 3
% Compile with: pdflatex warlouzet-hague-to-maastricht-reading-summary.tex

\documentclass[11pt,a4paper]{article}
\usepackage[utf8]{inputenc}
\usepackage[T1]{fontenc}
\usepackage[margin=2.5cm]{geometry}
\usepackage{booktabs}
\usepackage{enumitem}
\setlist{nosep,leftmargin=*}
\usepackage{parskip}
\usepackage{microtype}
\usepackage{hyperref}
\usepackage{xcolor}
\definecolor{eublue}{RGB}{0,51,153}

\hyphenation{Spitzenkandidat Spitzenkandidaten}
\hypersetup{
  colorlinks=true,
  linkcolor=eublue,
  urlcolor=eublue,
  pdftitle={Warlouzet – Hague to Maastricht},
  pdfauthor={Introduction to the EU, UCM}
}

\begin{document}

\title{From The Hague summit to the Maastricht Treaty, 1969--1992\\Chapter from \textit{Reinventing Europe}}
\author{Introduction to the European Union\\Universidad Complutense de Madrid\\Week 3}
\date{}
\maketitle
\vspace{-1em}
\hrule
\vspace{1em}

\section*{Rhetorical Framework (Ethos, Logos, Pathos)}

\begin{table}[h]
\centering
\begin{tabular}{@{}lll@{}}
\toprule
\textbf{Appeal} & \textbf{Meaning} & \textbf{Warlouzet's Use} \\
\midrule
\textbf{Ethos} & Credibility, authority, trust & Institutional history; synthesis of treaties and summits; scholarly narrative \\
\textbf{Logos} & Logic, structure, cause-and-effect & Chronological arc (crisis $\to$ sclerosis $\to$ renewal); institutional development \\
\textbf{Pathos} & Emotion, stakes, human dimension & Crisis, paralysis, renewal; oil shock, recession, revival of integration \\
\bottomrule
\end{tabular}
\end{table}

\textbf{Key insight:} Warlouzet structures the chapter around a \textit{logos} arc---turning point (Hague) $\to$ crisis (oil) $\to$ stagnation (Eurosclerosis) $\to$ renewal (SEA, Maastricht). \textit{Pathos} appears in the drama of crisis and recovery; \textit{ethos} in the authoritative treatment of treaties and summits.

\section*{Bibliographic Information}

\begin{itemize}
\item \textbf{Author:} Laurent Warlouzet
\item \textbf{Title:} From The Hague summit to the Maastricht Treaty, 1969--1992
\item \textbf{Source:} In LEUCHT, Brigitte; SEIDEL, Katja \& WARLOUZET, Laurent, \textit{Reinventing Europe}
\item \textbf{Publisher:} Bloomsbury, London
\item \textbf{Year:} 2023
\end{itemize}

\section*{Summary \& Key Points}

\begin{itemize}
\item The Hague Summit (1969) marked a turning point, setting agenda for enlargement and economic and monetary union.
\item First enlargement in 1973 (UK, Ireland, Denmark) expanded Community from 6 to 9 members.
\item The 1973 oil crisis revealed weaknesses in European coordination and led to economic recession.
\item Period of ``Eurosclerosis'' in 1970s--1980s with slow integration progress.
\item Direct elections to European Parliament began in 1979, increasing its legitimacy.
\item The Single European Act (1986) and Maastricht Treaty (1992) marked major steps forward in integration.
\item This period saw transition from European Community to European Union.
\end{itemize}

\section*{Main Arguments}

\subsection*{I. The Hague Summit as Turning Point}

After De Gaulle's resignation, The Hague Summit in December 1969 set a new agenda focusing on enlargement, economic and monetary union, and deeper cooperation. This marked the end of the paralysis period.

\subsection*{II. First Enlargement and Its Challenges}

The 1973 enlargement brought in UK, Ireland, and Denmark, but coincided with economic crisis. The enlargement process revealed the need for special arrangements (Commonwealth, fisheries) and showed how expansion complicated integration.

\subsection*{III. The 1973 Oil Crisis}

The oil crisis revealed the Community's weakness---inability to coordinate a common response. Each country acted individually, showing limits of integration and leading to economic recession and stagflation.

\subsection*{IV. From Eurosclerosis to Renewal}

The 1970s--1980s saw slow progress (``Eurosclerosis''), but also important developments like direct elections to Parliament (1979) and the Single European Act (1986), which revived integration.

\subsection*{V. Maastricht Treaty}

The Maastricht Treaty (1992) created the European Union and established the framework for Economic and Monetary Union, representing a major step forward in integration.

\section*{Context \& Analysis}

The chapter traces the arc from post-de Gaulle renewal (Hague Summit) through crisis (oil shock) and stagnation (Eurosclerosis) to institutional revival (SEA, Maastricht). This period established the framework for EMU and debates about deepening vs.\ widening.

\section*{Relevance to Course}

\begin{itemize}
\item \textbf{Ethos:} Essential for understanding first enlargement and transition to EU.
\item \textbf{Logos:} Oil crisis, Eurosclerosis, SEA, Maastricht---cause-and-effect, institutional logic.
\item \textbf{Pathos:} Crises and renewal---stakes of integration; how crises shape (and sometimes revive) the project.
\item \textbf{All three:} Context for EMU and debates about deepening vs.\ widening.
\end{itemize}

\section*{Discussion Questions}

\begin{itemize}
\item Causes and consequences of 1973 oil crisis? (Ethos: who responded? Logos: causal chain? Pathos: stakes, vulnerability)
\item How did first enlargement affect the Community? (Ethos: who gained, who lost? Logos: structural effects; Pathos: tension, complexity)
\item Why ``Eurosclerosis'' and what ended it? (Ethos: who diagnosed it? Logos: what broke the deadlock? Pathos: stagnation vs.\ renewal)
\item Significance of Maastricht Treaty? (Ethos: who framed it? Logos: institutional change; Pathos: achievement, culmination)
\end{itemize}

\vspace{1em}
\noindent\footnotesize\textit{Reference:} Warlouzet, Laurent. ``From The Hague summit to the Maastricht Treaty, 1969--1992.'' In LEUCHT, Brigitte; SEIDEL, Katja \& WARLOUZET, Laurent, \textit{Reinventing Europe}, London: Bloomsbury, 2023.

\end{document}
